
Our study connects research on LLM benchmarking, field experiments on AI-assisted work, and emerging frameworks for evaluating human--AI complementarity.

\subsection{LLM Benchmarks and Autonomous Performance}

A large computer-science literature evaluates LLMs by asking them to complete tasks autonomously. Early general-purpose benchmarks emphasize knowledge and reasoning across broad collections of tasks. MMLU, for example, tests knowledge across 57 subjects \citep{hendrycks2021}, while HELM evaluates models across a wide range of scenarios and evaluation dimensions \citep{liang2023}. These benchmarks have played an important role in standardizing model comparisons, but their basic unit of evaluation remains the output produced by a model working independently.

More recent benchmarks have moved toward realistic professional and economically grounded tasks. These evaluations test models on industrial workflows \citep{sopbench2025}, occupational tasks spanning specialized domains \citep{occubench2026}, and work products designed to approximate economically valuable professional activity \citep{gdpval2025, profbench2025}. JobBench further makes task selection more worker-centered by constructing 130 agentic tasks across 35 occupations from duties that workers report wanting to delegate, and by evaluating the resulting work against fact-anchored criteria \citep{jobbench2026}. This represents an important shift from selecting tasks solely according to economic value or technical feasibility toward also considering which activities workers want AI systems to perform. This progression has increased the realism and economic relevance of model evaluation. Nevertheless, operationally, these benchmarks still evaluate the model as the party responsible for completing the selected task and producing the final deliverable end-to-end. They therefore provide increasingly realistic measures of automation, but do not directly measure the value a model creates by augmentation.

That focus leaves a distinct capability undermeasured. In many workplace applications, the model does not directly produce the final deliverable. Instead, it helps another party plan, execute, review, or revise the work. Evidence that autonomous performance does not necessarily predict interactive performance reinforces this concern \citep{collins2024, chang2025}. A benchmark that identifies the best independent task solver may therefore provide limited guidance about which model should be deployed as an assistant.

\subsection{AI Augmentation and Workplace Productivity}

A parallel literature in economics and management examines how access to generative AI affects worker performance. Across several settings, experiments have found substantial productivity effects. ChatGPT reduced the time required for professional writing tasks while improving output quality \citep{noy2023}; GitHub Copilot accelerated completion of a programming task \citep{peng2023}; and a large-scale deployment of a generative-AI assistant increased productivity among customer-support workers \citep{brynjolfsson2025}. Related studies examine consultants \citep{dellacqua2023}, advertising professionals \citep{chen2024, ju2026}, and creative writers \citep{doshi2024}.

These studies establish that generative AI can affect worker productivity and that its effects depend on the task, worker, and mode of interaction. At the same time, field experiments are necessarily narrow by design. They typically study one model or tool, one population, and one domain, often relative to an unaided control condition. Their purpose is generally to estimate whether access to a particular AI system improves performance in a particular setting, rather than to compare which of several models would provide the greatest assistance to the same worker.

The evidence also cautions against treating human--AI collaboration as automatically beneficial. In a meta-analysis of 106 experiments, \citet{vaccaro2024} find that human--AI combinations frequently underperform the better of the human-alone and AI-alone conditions. The value of assistance may depend on how capabilities are distributed between the parties, how work is allocated, and whether the system's guidance is appropriately calibrated to the worker. These findings motivate evaluations that compare not only whether assistance is available, but also which model provides that assistance and how its value varies across tasks.

This question is economically consequential because automation and augmentation can produce different organizational and distributional effects. \citet{agrawal2023} argue that automating one worker's tasks can create augmentation opportunities elsewhere in a production process, while \citet{marguerit2025} distinguish the labor-market effects of augmentation-oriented and automation-oriented AI. Worker preferences may also favor collaborative arrangements: a nationwide survey of 1,500 U.S. workers across 104 occupations found that ``equal partnership'' with AI, rather than replacement, was the dominant preference in nearly half of the occupational categories studied \citep{shao2025}. Together, this research highlights the importance of evaluating AI systems in the roles through which they are actually deployed.

\subsection{Automation, Augmentation, and Human--AI Complementarity}

A growing body of work more directly distinguishes automation from augmentation. \citet{haupt2025} argue that evaluation by imitation privileges systems that reproduce human outputs and replace human effort, rather than systems designed to improve and augment human performance. They call for evaluations in which humans and AI systems solve tasks together. This critique shifts the object of evaluation from the model's standalone capability to the value it creates within a larger workflow.

\citet{fugener2026} formalize the distinction between automation and augmentation, showing that the optimal allocation of work depends on the structure of complementarity between human and AI capabilities. A system that performs well independently need not create the greatest value when paired with another decision-maker or worker. Their analysis therefore provides a theoretical basis for expecting model rankings to vary across autonomous and assistive roles. A growing set of efforts evaluates this type of assistance empirically, but each within a narrow scope. Upwork's Human+Agent Productivity Index (HAPI) evaluates frontier agents on approximately 300 paid client projects, and finds that expert involvement can substantially improve agent performance \citep{upwork_hapi2025}. However, a human expert assists the AI, reviewing and correcting the agent's output, so the human occupies the coaching role and the model is the partly being helped. MathTutorBench benchmarks the pedagogical quality of LLM tutors, finding that stronger problem-solving ability does not automatically imply better teaching \citep{macina2025}. \citet{gandhi2025} study strong--weak model collaboration for code generation, in which an expert model supplies plans for an inferior model to execute, reducing cost while preserving performance. \citet{saha2023} examine a teacher--student simulation between two LLMs, showing that the stronger model's explanations improve the weaker model's performance when explanations are personalized to the student. Each of these studies is confined to a single domain (e.g. tutoring, coding, freelance projects) and to a single human--LLM or model--model pairing, rather than comparing many models against one another in the assistant role.

Our approach differs from these studies by spanning a range of real-world professional tasks and placing models into an assistance/guidance role to a fixed worker. Importantly, the latter allows measurement of the marginal value of an assistant model's guidance rather than its joint outcome as a team. Rather than asking only whether automation or augmentation is optimal in a particular setting, we compare the same set of models under both modes and across multiple professional domains. This design allows us to examine whether relative model rankings are stable across roles and whether a model's autonomous capability is also an indication of the value of its guidance to a fixed downstream worker. In sum, our study operationalizes this distinction as a cross-model, cross-mode, and cross-task empirical evaluation.


%%%%% OLD %%%%%%%

%A growing body of research has attempted to quantify the economic impact of AI by measuring its performance on economically relevant tasks. Early work focused predominantly on automation, but recent scholarship has shifted toward augmentation—the idea that AI can amplify human capability rather than simply replace it. Representative examples include: work arguing that AI raises the returns to worker skill rather than eroding them (NBER, 2023); studies showing that LLM-powered tools enhance cognitive tasks in ways that traditional digital automation cannot (2025); findings that augmentation tends to favor high-skilled occupations while automation disproportionately displaces low-skilled ones, with implications for wage inequality (arXiv, 2025); and policy-oriented analyses distinguishing "simplistic automation" from genuine augmentation as a strategic choice (Berkeley Roundtable, 2024).

%Taken together, this literature establishes that model performance is multidimensional: it depends on which task is being evaluated, which model is being used, and whether the model is operating in automation or augmentation mode. Yet existing evaluations fail to capture this multidimensionality in a unified way. They either (1) treat AI as a monolithic entity, testing a single model as a stand-in for "AI as a whole"; (2) study augmentation in narrow field experiments that are hard to generalize; or (3) benchmark models broadly but almost exclusively in automation and one-shot settings, with no general framework that jointly addresses task selection, usage mode, and cross-model comparison. A systematic, cross-model evaluation that asks which models are best suited for which tasks and for which usage mode does not currently exist.

% min min - this is my suggested lit review writing. we can mesh them / incorporate better though.

%\textbf{Min Version - feel free to scrap or edit}

%The clearest evidence that generative AI has economic value comes from a set of controlled experiments on professional work. \citet{noy2023} show that ChatGPT reduced task completion time by 40\% and raised output quality by 18\% in professional writing. \citet{peng2023} find that GitHub Copilot cut programming time by more than half. \citet{brynjolfsson2025} document a 15\% productivity gain in customer support, with the largest gains accruing to less experienced workers. \citet{dellacqua2023} extend this to knowledge work, finding that GPT-4 made BCG consultants substantially more productive on tasks inside the model's capability frontier---but actively hurt performance on tasks outside it. Yet when \citet{vaccaro2024} synthesize 106 experimental studies on human-AI collaboration, they find that human-AI combinations on average perform \emph{worse} than the best of either alone, with the largest losses concentrated in decision-making tasks. The contradiction between these bodies of evidence is not noise. It is a signal that ``does AI help?'' is the wrong question. The answer depends critically on what role the AI is playing and on what kind of task.

%Several studies point toward usage mode as the key moderating variable. \citet{fugener2026} provide both a theoretical framework and empirical validation distinguishing automation, in which AI substitutes for human work, from augmentation, in which AI supports human judgment. They show that the optimal mode depends on the type of complementarity between human and AI capability. \citet{chen2024} find that LLMs used as sounding boards improve output quality for non-experts, while LLMs used as ghostwriters can harm expert performance through anchoring effects. \citet{doshi2024} show that generative AI raises individual creativity in short-story writing but reduces collective diversity by making outputs more similar. \citet{ju2026} find that human-AI teams produce more output and higher-quality text than human-human teams, but lower-quality images and more homogeneous outputs overall. Taken together, these studies establish that the automation-augmentation distinction is empirically real and consequential. But each tests a single model on a single task domain, making it impossible to ask which models are better suited to which mode, or whether the pattern holds across professionally diverse work.

%The standard infrastructure for comparing models across tasks offers little traction on this question. \citet{hendrycks2021} introduced MMLU as a broad multitask benchmark covering 57 subjects; \citet{liang2023} extended this with HELM, which evaluates 30 models across 42 scenarios using seven metrics. \citet{srivastava2023} contributed BIG-Bench, demonstrating that model performance varies substantially across 204 tasks. These benchmarks improve breadth and standardization, but they evaluate models as isolated solvers. The interactive dimension (how a model performs when a human is in the loop) is absent by design. \citet{collins2024} and \citet{chang2025} begin to address this gap: Collins et al. show that static evaluation diverges from interactive usefulness in mathematical theorem proving, and Chang et al. find that AI-alone accuracy fails to predict user-AI accuracy on conversational reformulations of MMLU. Both papers, however, only work with one or two models and academic task content. \citet{gdpval2025} take a step toward economically grounded evaluation by benchmarking models on tasks drawn from 44 occupations, using expert-created deliverables and pairwise comparison---but the framework evaluates models only as end-to-end task performers, with no augmentation regime.

%What the literature does not yet provide is a framework that jointly varies usage mode and task domain across multiple models in a single unified pipeline. \citet{fugener2026} come closest, but their empirical study covers a single task type and does not produce cross-model rankings. The result is that practitioners and researchers lack the comparative map needed to ask: which models are best suited to replacing work, which to improving work, and does the answer change across professionally meaningful domains? We propose a framework that addresses this.